% მონახაზი: თავი 1 (შესავალი/მოტივაცია) და თავი 3 (მოთხოვნათა ანალიზი)
% ბაკალავრის ნაშრომი ინფორმატიკაში, ბრემენის უნივერსიტეტი — სერგი კონიაშვილი
% თარიღი: 2026-11-01 — სადისკუსიო საფუძველი ხელმძღვანელებთან შეხვედრისთვის (2026-11-02).
% ფაქტები ეფუძნება ექსკლუზიურად docs/thesis/THESIS_KNOWLEDGE_BASE.md-ს (VERIFIED).
% ჯერ კიდევ დაუმოწმებელი ადგილები მონიშნულია [TODO]-თი.
\documentclass[11pt,a4paper]{article}
\usepackage{fontspec}
\usepackage[a4paper,margin=2.5cm]{geometry}
\usepackage{booktabs}
\usepackage{longtable}
\usepackage{array}
\usepackage{xcolor}
\usepackage{enumitem}
\usepackage[hidelinks]{hyperref}

\setmainfont{Sylfaen}

\newcommand{\enquote}[1]{,,#1``}
\newcommand{\todo}[1]{\textcolor{red!70!black}{[TODO: #1]}}
\newcommand{\kitkhva}[1]{\par\smallskip\noindent\colorbox{blue!8}{\parbox{\dimexpr\linewidth-2\fboxsep}{\small\textbf{კითხვა ხელმძღვანელებს:} #1}}\par\smallskip}

\title{მონახაზი: მოტივაცია და ლიტერატურაზე დაფუძნებული მოთხოვნათა ანალიზი\\[0.3em]
\large კლიმატისა და ჯანდაცვის მონაცემთა LLM-ზე დაფუძნებული ინტეგრაცია}
\author{სერგი კონიაშვილი\\ ბაკალავრიატი ინფორმატიკაში, ბრემენის უნივერსიტეტი}
\date{მდგომარეობა: 1 ნოემბერი, 2026 (განახლებული) -- სადისკუსიო საფუძველი}

\begin{document}
\maketitle

\noindent\textit{შენიშვნა:} ეს მონახაზი წარმოადგენს ნაშრომის დაგეგმილ სტრუქტურას,
ასევე პირველი და მესამე თავების წინასწარ სამუშაო ვერსიებს. წითლად მონიშნული
ადგილები ჯერ არ არის სრულად დამოწმებული ან ღიაა. ციტირების სტილი დროებით
არის ავტორი-წელი (იხ. კითხვა სექციაში~\ref{sec:questions_ka}).

\section*{ნაშრომის დაგეგმილი სტრუქტურა}
\begin{enumerate}[nosep]
  \item შესავალი (მოტივაცია, პრობლემის დასმა, მიზანი, სტრუქტურა)
  \item საფუძვლები (კლიმატი და ვექტორული დაავადებები; მონაცემთა წყაროები:
        შემთხვევები, რეანალიზის კლიმატური მონაცემები, მოსახლეობა; ETL; LLM)
  \item მოთხოვნათა ანალიზი (ნაწარმოები სანიმუშო პუბლიკაციებიდან)
  \item მონათესავე ნაშრომები და არსებული ინსტრუმენტები \todo{ინსტრუმენტების შემოწმება ოფიციალური დოკუმენტაციიდან}
  \item კონცეფცია და არქიტექტურა
  \item იმპლემენტაცია
  \item ევალუაცია (საქმის შესწავლა CS1--CS3)
  \item დისკუსია, შეზღუდვები და სამომავლო პერსპექტივები
  \item დასკვნა
\end{enumerate}

% ---------------------------------------------------------------------------
\section{შესავალი}

\subsection{მოტივაცია}

კლიმატის ცვალებადობასა და ვექტორულ დაავადებებს (მაგ. დენგეს ცხელება) შორის
ურთიერთქმედება სულ უფრო ხშირად შეისწავლება მონაცემებზე დაფუძნებული მოდელებით.
აქტუალური მაგალითია Suttirat et al.\ (2026)-ის კვლევა, რომელიც დენგეს ჰემორაგიული
ცხელების (DHF) ყოველთვიურ შემთხვევებს ტაილანდის 77 პროვინციაში 20 წლის მანძილზე
(2003--2022) აკავშირებს კლიმატურ, დისტანციური ზონდირებისა და სოციო-ეკონომიკურ
მონაცემებთან და აანალიზებს XGBoost-ისა და SHAP-ის მეშვეობით (გვ.~5--8). Lowe et al.\
(2021) მსგავსად იკვლევენ ჰიდრომეტეოროლოგიური ექსტრემუმებისა და ურბანიზაციის
გავლენას დენგეს რისკზე ბრაზილიის 558 მიკრორეგიონში (გვ.~e211--e212).

ორივე კვლევას აერთიანებს ის გარემოება, რომ უშუალო სტატისტიკურ და მანქანური
სწავლების მოდელირებას წინ უსწრებს მონაცემთა მომზადების ფართო, მრავალსაფეხურიანი
მილსადენი. Suttirat et al.\ ამ წინასწარ ნაბიჯებს შემდეგნაირად აღწერენ:

\begin{quote}
\enquote{Gridded climatic and remote-sensing data were spatially aggregated to
the provincial level; lagged variables and bioclimatic indices were then derived
from monthly provincial temperature and precipitation series. Socioeconomic data
were obtained directly at the provincial level: the biennially reported household
income series was linearly interpolated to fill missing years, and all
socioeconomic series were then upsampled to monthly resolution by assigning each
year's reported value (typically released at year-end) to every month of that
year.} (Suttirat et al., 2026, გვ.~4)
\end{quote}

მონაცემთა დამუშავების ეს პროცედურები არ არის რომელიმე ცალკეული სტატიის
სპეციფიკური თვისება. Lowe et al.\ (2021) თითქმის იდენტურ ნაბიჯებს ახორციელებენ:
ისინი იღებენ ყოველთვიურ შემთხვევებს ეროვნული შეტყობინების სისტემიდან,
მოსახლეობის წლიური შეფასებებით ითვლიან ინციდენტობას 100\,000 მცხოვრებზე,
ახდენენ გრიდირებული კლიმატური მონაცემების (Tmin, Tmax, PDSI) ფართობით შეწონილ
აგრეგაციას ანალიზურ ერთეულებზე და კლიმატის მონაცემებს იღებენ დამატებითი
ერთი წლით ადრე პირველ სადამკვირვებლო შემთხვევამდე, \enquote{to allow for a
lag period before the first dengue observation in January, 2001} (გვ.~e211).
Suttirat et al.\ ასევე აგენერირებენ დროში დაყოვნებულ ცვლადებს (1--3 თვიანი
Lags) ყველა კლიმატური და გარემოსდაცვითი მახასიათებლისთვის (გვ.~7).
ცხრილი~\ref{tab:steps_ka} ერთმანეთს ადარებს ორივე კვლევის მონაცემთა მომზადების ნაბიჯებს.

უფრო მეტიც, მიუხედავად იმისა, რომ პირველადი მონაცემები ხშირად საჯაროდ
ხელმისაწვდომია, ისინი უშუალოდ ანალიზისთვის გამზადებული არ არის. Lowe et al.\
სრულად იყენებენ ღია მონაცემებს (გვ.~e218), თუმცა ხუთი ჰეტეროგენული წყაროდან,
რომელთაც აქვთ განსხვავებული ფორმატები და დროითი რეზოლუციები (თვიური, წლიური,
სტატიკური, 0.5°-იანი გრიდი); მონაცემთა ბაზის თავმოყრა თანაავტორთა ცალკე
დავალებად არის მითითებული (გვ.~e218). ჯანდაცვის საზედამხედველო მონაცემებთან
დაკავშირებით Clarke et al.\ (2024) OpenDengue-ს ბაზის აღწერისას აღნიშნავენ:

\begin{quote}
\enquote{Dengue case data exists in multiple formats from a wide variety of
sources that require various processing methodologies. [\ldots] Locating,
extracting, processing and standardising data all takes time but is essential to
enable reuse and re-analysis.} (Clarke et al., 2024, გვ.~2)
\end{quote}

OpenDengue წარმატებით აგვარებს შემთხვევათა მონაცემების გლობალურ კონსოლიდაციას.
თუმცა, ამ მონაცემების გაერთიანება მეტეოროლოგიურ რეანალიზთან — რაც თავად Clarke et al.-ის
მიერ დასახელებულია, როგორც ერთ-ერთი უმთავრესი სამომავლო გამოყენება (გვ.~1, 9) —
მკვლევრის ინდივიდუალურ პასუხისმგებლობად რჩება.

დაბოლოს, მონაცემთა ინჟინერიის ამ მილსადენის შესრულება მოითხოვს პროგრამირების
ღრმა ცოდნას: Suttirat et al.-მა მონაცემთა მომზადება განახორციელეს Python-ში (pandas),
ხოლო ბიოკლიმატური ინდექსები გამოთვალეს R-ის პაკეტით \texttt{dismo} (გვ.~5, 8).
დარგობრივი სპეციალისტებისთვის, რომელთაც არ აქვთ პროგრამული ინჟინერიის გამოცდილება
(მაგ. საზოგადოებრივი ჯანდაცვისა და ეპიდემიოლოგიის მუშაკები), ეს ტექნიკური ბარიერი
სერიოზულ დაბრკოლებას ქმნის \todo{წყარო Public Health-ში პროგრამირების/Data Science უნარების დეფიციტზე}.
არსებული პროგრამული ხელსაწყოები კი ამ ნაბიჯებს მხოლოდ ფრაგმენტულად ფარავენ
\todo{მონათესავე ინსტრუმენტები: Copernicus CDS, Google Earth Engine, DHIS2 Climate \& Health, IRI ENACTS -- ოფიციალური დოკუმენტაციიდან მოწმობა}.

\subsection{პრობლემის დასმა და მიზნები}

წინამდებარე ნაშრომის მიზანია ისეთი პროგრამული პლატფორმის შექმნა, რომელიც
მკვლევრებს სრულად მოუხსნის მოდელირების წინამორბედ მონაცემთა მომზადების
ტვირთს — წყაროების მოძიებას, გაწმენდას, ჰარმონიზაციას, ინტეგრაციასა და
დოკუმენტირებას — პროგრამირების ცოდნის გარეშე, სრული გამჭვირვალობისა და
აუდიტირებადობის პირობებში, რათა მკვლევრებს შეეძლოთ მთელი ყურადღება
დაუთმონ უშუალო სამეცნიერო მოდელს. მომხმარებელი ინტერაქტიული ფორმით ან
ბუნებრივი ენით აღწერს საჭირო მონაცემთა ნაკრებს (დაავადება, რეგიონი, პერიოდი,
კლიმატის ცვლადები); Large Language Model (LLM) ამ მოთხოვნას გარდაქმნის
სტრუქტურირებულ მოთხოვნად. სისტემა ირჩევს შესაფერის ავტორიტეტულ
წყაროებს, ჰარმონიზაციას უკეთებს რეზოლუციებს, აერთიანებს ჩანაწერებს ერთ
ანალიზისთვის გამზადებულ ცხრილში, ითვლის ნაწარმოებ სიდიდეებს (ინციდენტობა,
დაყოვნებული ცვლადები) და აღრიცხავს ტრანსფორმაციის თითოეულ ნაბიჯს.

LLM-ის როლი განზრახ შეზღუდულია. ის ბუნებრივი ენის მოთხოვნას ფორმის ველებად
გარდაქმნის, რომელთა შემოწმება და შეცვლა გაშვებამდე შეიძლება, შედეგს გასაგები
ენით ხსნის და ჩატ-ასისტენტში ეხმარება დარგობრივ კითხვებსა და ფორმის
შევსების შეთავაზებებში. მონაცემების მოპოვება, გაწმენდა, ჰარმონიზაცია,
გაერთიანება, გამოთვლები და ექსპორტი კი მხოლოდ დეტერმინისტული პროგრამული
კომპონენტებით სრულდება; ამიტომ ერთი და იგივე მოთხოვნა ერთსა და იმავე ცხრილს
იძლევა. LLM-ის გამოსავალი მაინც საჭიროებს შემოწმებას: საკუთარ ტესტში მოდელმა
არამნიშვნელოვანი კორელაცია (lag~3, $n = 9$) ყველაზე ძლიერ კავშირად დაასახელა.
ამის შემდეგ ის ყოველ კორელაციასთან ერთად იღებს მნიშვნელოვნების ნიშანს და
წყვილების რაოდენობას და დასკვნად მხოლოდ მნიშვნელოვანი კავშირების დასახელება
შეუძლია. რამდენად საიმედოდ მუშაობს LLM-ის კომპონენტები, ეს ევალუაციაში
შემოწმდება \todo{კრიტერიუმების განსაზღვრა, მაგ.\ სწორად შევსებული ფორმის ველების წილი}.

ხაზგასმით უნდა აღინიშნოს, რომ სისტემა ცალსახად \emph{არ წარმოადგენს} საპროგნოზო
მოდელს: ის ამზადებს სტანდარტიზებულ მონაცემებს შემდგომი ანალიტიკური მოდელებისთვის.
ამცირებს თუ არა სისტემა რეალურად ხელით შრომას, ეს ემპირიულად გაიზომება
ევალუაციაში (თავი~7, საქმის შესწავლა CS1: ერთი და იმავე სამიზნე მონაცემთა
ნაკრების ხელით აწყობა Python/pandas-ით vs. სისტემის გამოყენებით; დროის,
ხელით გადადგმული ნაბიჯებისა და შედეგების თანხვედრის შედარება).

\subsection{ნაშრომის სტრუქტურა}
\todo{დასრულდეს ნაშრომის საბოლოო სარჩევის დამტკიცების შემდეგ}

% ---------------------------------------------------------------------------
\setcounter{section}{2}
\section{მოთხოვნათა ანალიზი}

\subsection{მეთოდოლოგიური მიდგომა}

ფუნქციონალური და არაფუნქციონალური მოთხოვნები სისტემურად არის გამოყვანილი
კლიმატისა და ჯანდაცვის სფეროს რეალური რეცენზირებული სამეცნიერო სტატიებიდან.
თითოეული კვლევის მეთოდების ნაწილიდან ექსტრაქტირებულია მოდელირებამდე ჩატარებული
მონაცემთა მომზადების ყველა ნაბიჯი. ყოველი ნაბიჯი, რომელიც აუცილებელია
ანალიზისთვის გამზადებული მონაცემთა ნაკრების მისაღებად, გარდაქმნილია სისტემის
კონკრეტულ მოთხოვნად ან დასაბუთებულად არის გამორიცხული (საზღვრების გავლება).
მიმართება: მომზადების ნაბიჯი $\rightarrow$ პუბლიკაცია (გვერდის მითითებით)
$\rightarrow$ სისტემის მოთხოვნა ასახულია Traceability-Matrix-ში (ცხრილი~\ref{tab:trace_ka}).

ამჟამად ემპირიულ საფუძველს წარმოადგენს სამი პუბლიკაცია:
\begin{itemize}[nosep]
  \item \textbf{Suttirat et al.\ (2026)}: დენგე, ტაილანდი, 77 პროვინცია,
        2003--2022, XGBoost/SHAP -- ძირითადი ვალიდაციური საქმის შესწავლა;
  \item \textbf{Lowe et al.\ (2021)}: დენგე, ბრაზილია, 558 მიკრორეგიონი,
        2001--2019, ბაიესური სივრცე-დროითი მოდელი -- ადასტურებს მონაცემთა
        მომზადების ნაბიჯების ტიპურობას;
  \item \textbf{Clarke et al.\ (2024)}: OpenDengue მონაცემთა ბაზა (Data Descriptor) --
        განსაზღვრავს მოთხოვნებს წარმომავლობაზე, ვერსიონირებასა და ეპიდზედამხედველობის
        მონაცემთა შედარებადობაზე.
\end{itemize}

\subsection{მონაცემთა მომზადების ნაბიჯები სანიმუშო პუბლიკაციებში}

\begin{longtable}{@{}p{0.27\linewidth}p{0.33\linewidth}p{0.33\linewidth}@{}}
\caption{მონაცემთა მომზადება მოდელირებამდე (ამონაწერი)}\label{tab:steps_ka}\\
\toprule
ნაბიჯი & Suttirat et al.\ (2026) & Lowe et al.\ (2021) \\
\midrule
\endfirsthead
\toprule
ნაბიჯი & Suttirat et al.\ (2026) & Lowe et al.\ (2021) \\
\midrule
\endhead
შემთხვევათა მოძიება & DHF-ის ყოველთვიური შემთხვევები, ეპიდემიოლოგიის ბიურო (გვ.~5) &
  ყოველთვიური შეტყობინებები, SINAN/DATASUS (გვ.~e211) \\
ადმინისტრაციული ერთეულები & Bueng Kan შეერწყა Nong Khai-ს (გვ.~5) & -- \\
კლიმატის ცვლადები & TerraClimate: ნალექი, Tmin, Tmax და სხვ. (გვ.~5) &
  CRU TS 4.04: Tmin, Tmax, PDSI (გვ.~e211) \\
სივრცითი აგრეგაცია & გრიდი პროვინციებზე საშუალოვდება, მეთოდი მითითებული არაა (გვ.~5) &
  ფართობით შეწონილი საშუალო \texttt{exactextractr}-ით (გვ.~e211) \\
დაყოვნებული ცვლადები & Lags 1--3 თვე (გვ.~7) & Lags 0--6 თვე (DLNM, გვ.~e212) \\
ბუფერი Lag-ისთვის & -- & კლიმატის მონაცემები შემთხვევებზე 1 წლით ადრე (გვ.~e211) \\
ინციდენტობა 100\,000-ზე & დიახ (გვ.~7) & დიახ, მოსახლეობის წლიური შეფასებით (გვ.~e211) \\
წლიური $\rightarrow$ თვიური & წლიური მნიშვნელობა მიეწერება ყველა თვეს (გვ.~4--5) &
  არ არის აღწერილი \\
გამოტოვებული წლები & შინამეურნეობების შემოსავალი წრფივად ინტერპოლირდება (გვ.~5) & -- \\
ნაწარმოები ინდექსები & 19 ბიოკლიმატური ინდიკატორი (\texttt{dismo}, გვ.~5) & -- \\
ხელსაწყოები & Python/pandas, R (გვ.~5, 8) & R, \texttt{exactextractr}, R-INLA (გვ.~e211--e212) \\
\bottomrule
\end{longtable}

\subsection{ნაწარმოები მოთხოვნები (Traceability-Matrix)}

\begin{longtable}{@{}p{0.25\linewidth}p{0.19\linewidth}p{0.40\linewidth}p{0.12\linewidth}@{}}
\caption{მიმართება: მომზადების ნაბიჯი $\rightarrow$ პუბლიკაცია $\rightarrow$ მოთხოვნა}\label{tab:trace_ka}\\
\toprule
ნაბიჯი / დაკვირვება & წყარო & მოთხოვნა & სტატუსი \\
\midrule
\endfirsthead
\toprule
ნაბიჯი / დაკვირვება & წყარო & მოთხოვნა & სტატუსი \\
\midrule
\endhead
ინციდენტობა 100\,000-ზე როგორც მიზანი & Suttirat გვ.~7; Lowe გვ.~e211 &
  FR-17: მოსახლეობის მონაცემების მიბმა და ინციდენტობის გამოთვლა 100\,000-ზე & დანერგილია \\
წლიური მოსახლეობა, თვიური შემთხვევები & Suttirat გვ.~4--5 (წლიური მნიშვნელობა ყველა თვეს ენიჭება); Lowe: მეთოდი აღწერილი არ არის &
  FR-17: წლიური მნიშვნელობების წრფივი ინტერპოლაცია თვეებზე -- \emph{საკუთარი გადაწყვეტილება}, რომ ინციდენტობა წლების საზღვარზე არ „ხტებოდეს“; განსხვავება Suttirat-ის მეთოდისგან მე-8 თავში განიხილება & დანერგილია \\
დაყოვნებული კლიმატის ცვლადები & Suttirat გვ.~7 (1--3); Lowe გვ.~e212 (0--6) &
  FR-23: დაყოვნებული მახასიათებლების სვეტები 1--6 მონაცემთა ნაკრებში & დანერგილია \\
კლიმატის მონაცემები შემთხვევების დაწყებამდე & Lowe გვ.~e211 &
  FR-23: Lag-ბუფერი საწყისი პერიოდების დანაკარგის თავიდან ასაცილებლად & დანერგილია \\
Tmin/Tmax მხოლოდ საშუალოს ნაცვლად & Suttirat გვ.~5; Lowe გვ.~e211 &
  FR-24: Tmin/Tmax დღიური ექსტრემუმების საშუალოს სახით & დანერგილია \\
განსხვავებული დროითი რეზოლუციები & Suttirat გვ.~4--5 &
  FR-20: რეზოლუციის ჰარმონიზაცია (მარაგისა და ნაკადის სიდიდეები, იხ. ქვემოთ) & დანერგილია \\
წარმომავლობა უნდა იყოს მიკვლევადი; ვერსია & Clarke გვ.~2, 4, 12 &
  FR-25: ექსპორტის მეტამონაცემები (წყაროები, ლიცენზიები, ვერსია, ნაბიჯები) & დანერგილია \\
დაფარვა ქვეყნებს შორის ცვალებადია & Clarke გვ.~1, 8 &
  FR-26: დროითი ხარვეზების ამოცნობა მიზეზებით, ავტომატური შევსების გარეშე & დანერგილია \\
წყაროებს შორის შედარებადობის შეფასება & Clarke გვ.~2 &
  FR-27: ორი წყაროს ემპირიული შედარება საერთო პერიოდში & დანერგილია \\
კავშირების ექსპლორაციული შემოწმება & Suttirat გვ.~7; Lowe გვ.~e212 &
  FR-19: კროს-კორელაციები Lags 0--3 (Pearson, Spearman) & დანერგილია \\
\bottomrule
\end{longtable}

\noindent FR-19-თან მიმართებით აუცილებელია მეთოდოლოგიური საზღვრის მკაფიოდ დაფიქსირება:
წრფივი კორელაციები ვერ იჭერს არაწრფივ ზღვრულ დინამიკას, როგორიც Suttirat et al.-მა
აღმოაჩინეს (დაახლოებით 21\,°C Tmin-ისთვის Lag~1-ით და 32\,°C Tmax-ისთვის Lag~3-ით, გვ.~12);
Lowe et al.\ ხაზს უსვამენ, რომ კლიმატის გავლენა გადაცემაზე არის \enquote{complex, non-linear,
and often delayed} (გვ.~e211). შესაბამისად, კორელაციის კოეფიციენტები ემსახურება
მხოლოდ მონაცემთა წინასწარ ექსპლორაციულ შემოწმებას და არა საბოლოო მოდელირებას.

\noindent დროითი რეზოლუციების ჰარმონიზაციისას (FR-17, FR-20) ორი სახის სიდიდე
უნდა გაიმიჯნოს. \emph{მარაგის სიდიდეები} (Bestandsgrößen) მდგომარეობას აღწერს
(მაგ.\ მოსახლეობა, შინამეურნეობის შემოსავალი); უფრო წვრილ რეზოლუციაზე მათი
მნიშვნელობა იგივე რჩება ან ცნობილ მნიშვნელობებს შორის ინტერპოლირდება. Suttirat
et al.\ წლიურ სოციალურ-ეკონომიკურ მნიშვნელობებს წლის ყველა თვეს ანიჭებენ
(გვ.~4--5), შემოსავლის გამოტოვებულ წლებს კი წრფივი ინტერპოლაციით ავსებენ
(გვ.~5); სისტემა წლიურ მოსახლეობას თვეებზე წრფივად ინტერპოლირებს (FR-17).
\emph{ნაკადის სიდიდეები} (Stromgrößen) კი დროში ჯამდება (მაგ.\ შემთხვევები,
ნალექი): გამსხვილებისას ჯამდება, დაწვრილებისას კი ისე ნაწილდება, რომ ჯამი
შენარჩუნდეს (FR-20, NFR-06). თვიური შემთხვევების რიცხვი მარაგის სიდიდესავით
თვის ყოველ დღეს რომ მივანიჭოთ, შემთხვევები დაახლოებით ოცდაათჯერ გაიზრდებოდა.
საშუალო სიდიდეები, მაგ.\ ტემპერატურა, გამსხვილებისას საშუალოდება.

\subsection{მოთხოვნები პირდაპირი ლიტერატურული წყაროს გარეშე}

მოთხოვნათა საწყისი კატალოგი (FR-01-დან FR-22-მდე) შეიქმნა ლიტერატურის ღრმა ანალიზამდე.
მისი ნაწილი პირდაპირ მტკიცდება ლიტერატურით (FR-17, FR-19, FR-20, იხ. ზემოთ). სხვა
მოთხოვნები კი გამომდინარეობს მიზნიდან \enquote{პროგრამირების ცოდნის გარეშე} (ბუნებრივი ენის
ინტერპრეტაცია FR-01, გასაგები ტექსტური ახსნა FR-09, Copilot ასისტენტი FR-13, მსოფლიო რუკა FR-22,
PDF ანგარიში FR-21) და ხელმძღვანელების მითითებებიდან (მაგ. FR-03--FR-08).
\kitkhva{უნდა დასაბუთდეს ეს მოთხოვნები ძირითადად Usability / Zero-Programming მიზნებით, თუ მიზანშეწონილია მათი შემცირება საინჟინრო კატალოგში?}

\subsection{ნაშრომის ფარგლები (Abgrenzung)}

ლიტერატურაში დოკუმენტირებული მონაცემთა მომზადების შემდეგი ნაბიჯები შეგნებულად არ
არის რეალიზებული; ისინი მე-8 თავში შეზღუდვებად განიხილება:
\begin{itemize}[nosep]
  \item სუბნაციონალური ერთეულები და გრიდირებული რასტრის სივრცითი საშუალო
        (Suttirat გვ.~5; Lowe გვ.~e211) -- პლატფორმა ამჟამად იყენებს ერთ წარმომადგენლობით
        კოორდინატს თითოეული ქვეყნისთვის;
  \item დროთა განმავლობაში ცვალებადი ადმინისტრაციული საზღვრების ჰარმონიზაცია (Suttirat გვ.~5);
  \item ბიოკლიმატური ინდიკატორები (მაგ. WorldClim bio1--bio19) და სატელიტური ზონდირება (Suttirat გვ.~5);
  \item ხარვეზების ავტომატური შევსება სხვა წყაროდან: სხვადასხვა საზედამხედველო სისტემის
        მონაცემები განსხვავებულ განსაზღვრებებს ეყრდნობა — OpenDengue ითვლის ყველა საეჭვო
        და დადასტურებულ შემთხვევას (Clarke გვ.~5), მაშინ როცა Suttirat et al.\ მხოლოდ DHF-ზე
        ფოკუსირდებიან (გვ.~5) —, რის გამოც გაუაზრებელი შერევა დაარღვევდა შედარებადობას (A20).
        პლატფორმა ამოიცნობს ხარვეზებს (FR-26) და სთავაზობს წყაროების მკაფიო, არა-დესტრუქციულ
        შედარებას (FR-27).
\end{itemize}

% ---------------------------------------------------------------------------
\section*{ღია კითხვები ხელმძღვანელებთან შეხვედრისთვის}\label{sec:questions_ka}
\begin{enumerate}[nosep]
  \item შეესაბამება თუ არა მოტივაციის არგუმენტაციის ჯაჭვი (მოდელები $\rightarrow$
        შრომატევადი, განმეორებადი მონაცემთა მომზადება $\rightarrow$ პროგრამირების ბარიერი
        $\rightarrow$ ავტომატიზებული პლატფორმა) თქვენს ხედვას?
  \item საკმარისია თუ არა ორი ემპირიული კვლევა და ერთი მონაცემთა ბაზის აღწერა,
        თუ სასურველია დამატებითი დაავადებების (მაგ. მალარია) ჩართვა?
  \item მიიჩნევა თუ არა Traceability-Matrix შესაფერის ფორმად ინფორმატიკის ბაკალავრის ნაშრომში მოთხოვნების დასაბუთებისთვის?
  \item კვლევითი კითხვის შეთავაზება: \enquote{რამდენად შეუძლია LLM-ზე დაფუძნებულ
        სისტემას კლიმატისა და ჯანდაცვის ჰეტეროგენული მონაცემების ინტეგრაციის
        ავტომატიზაცია ისე, რომ მონაცემთა მომზადების სისწორე, განმეორებადობა და
        მიკვლევადობა შემოწმებადი გახდეს?} შეესაბამება ეს მიმართულება?
  \item ციტირების რომელი სტილია სასურველი (ავტორი-წელი თუ ნუმერაციული)?
  \item ნაშრომის სტრუქტურა: შეესაბამება თუ არა მოთხოვნათა ანალიზის შემდეგ სექცია \enquote{მონათესავე ნაშრომების} განთავსება თქვენს მოლოდინს?
\end{enumerate}

\section*{გამოყენებული ლიტერატურა (მონახაზი)}
\begin{small}
\begin{description}[nosep,leftmargin=1em,labelindent=0pt]
\item Clarke, J., Lim, A., Gupte, P., Pigott, D.\,M., van Panhuis, W.\,G., \&
  Brady, O.\,J. (2024). A global dataset of publicly available dengue case count
  data. \emph{Scientific Data}, 11, 296.
  \url{https://doi.org/10.1038/s41597-024-03120-7}
\item Lowe, R., Lee, S.\,A., O'Reilly, K.\,M., Brady, O.\,J., Bastos, L.,
  Carrasco-Escobar, G., et al.\ (2021). Combined effects of hydrometeorological
  hazards and urbanisation on dengue risk in Brazil: a spatiotemporal modelling
  study. \emph{The Lancet Planetary Health}, 5(4), e209--e219.
  \url{https://doi.org/10.1016/S2542-5196(20)30292-8}
\item Suttirat, P., Chadsuthi, S., Aekthong, S., Rocklöv, J., Bicout, D.\,J.,
  Haddawy, P., Su Yin, M., Lawpoolsri, S., \& Modchang, C. (2026). Machine
  learning reveals temperature as a key predictor of dengue risk across
  Thailand's provinces: A 20-year analysis. \emph{PLoS Neglected Tropical
  Diseases}, 20(7), e0014590.
  \url{https://doi.org/10.1371/journal.pntd.0014590}
\end{description}
\end{small}

\end{document}
